\section*{Aktivitas Pembelajaran dan Praktikum}
\addcontentsline{toc}{section}{Aktivitas Pembelajaran dan Praktikum}

\begin{table}[H]
\centering
\caption{Rencana aktivitas Bab II}
\label{tab:ch02-aktivitas}
\begin{tabularx}{\textwidth}{@{}llX@{}}
\toprule
\textbf{Sesi} & \textbf{Bentuk} & \textbf{Kegiatan} \\
\midrule
1 & Ceramah & Himpunan, relasi, fungsi, dan induksi \\
2 & Latihan & Membuktikan sifat string dan bahasa \\
3 & Lokakarya & Menulis tata bahasa dan derivasi untuk bahasa sederhana \\
4 & Diskusi & Mengklasifikasikan bahasa pada hierarki Chomsky \\
\bottomrule
\end{tabularx}
\end{table}

\subsection*{Lokakarya Notasi}
\begin{enumerate}[label=\arabic*.,leftmargin=*]
  \item Diberikan bahasa \(L = \{w \in \{0,1\}^{*} \mid w \text{
        berakhir dengan } 01\}\), tuliskan lima string anggota dan
        lima string bukan anggota.
  \item Tuliskan tata bahasa bebas konteks untuk \(L\).
  \item Buat derivasi leftmost untuk string \(00101\).
  \item Gambarkan pohon urai untuk derivasi tersebut.
  \item Tentukan kelas bahasa \(L\) pada hierarki Chomsky dan
        jelaskan alasannya.
\end{enumerate}

Bab ini tidak menambahkan increment proyek. Milestone M0 dimulai pada
Bab III setelah arsitektur kompilator dibahas. Namun, mahasiswa
disarankan menyimpan catatan notasi dan contoh derivasi sebagai
rujukan saat menulis spesifikasi Subset-C.
